\section{RISK MODELS TRAINED ON FINDING HISTORIES}\label{sec:scoring}

Audit organizations increasingly build models that score programs for risk and
validate them the way any predictive model is validated: fit on one period,
test on another, report an AUC, check calibration. The metrics are not broken.
They answer a precise question about the label they are given, and the label
is the recorded finding. The error is to carry the result over and describe
the model as predicting where problems are. It has teeth because a score that
steers where auditors look can make itself right: a flagged program gets more
attention, more attention raises detection, higher detection produces recorded
findings, and the recorded findings validate the score
\citep{ensign2018,lum2016,lakkaraju2017,kleinberg2018}. The obvious objection
is that this is just a story about auditing more, so the construction worth
making holds average detection fixed.

Figure~\ref{fig:budget} does that. Across the whole horizontal axis the
average chance of catching an existing condition stays at
\mBudgetDetection{} and the overall recorded finding rate stays at
\mBudgetPrevalence{}. The only thing that changes is allocation: detection
moves away from the programs the score does not flag and toward the programs
it does. The score is statistically independent of where the conditions
actually are; it knows nothing. Measured against the recorded finding, its
AUC climbs from \mBudgetAucFstart{}, a coin flip, to \mBudgetAucFend{}, and
its Brier score improves. Measured against the underlying condition, its AUC
is exactly \mBudgetAucC{} at every point and its Brier score worsens by exactly
the amount the recorded-label Brier score improves.
On a validation report this reads as a model that became genuinely
discriminating, with no change in average detection. The calculation is
exact, from a single prespecified configuration.
Appendix~\ref{app:pred} gives the grid, the general result, a fitted rather
than an oracle score, and the finite-sample distributions at portfolio sizes
an audit organization actually has. Calibration is no stronger a guarantee: a
model can be perfectly calibrated for recorded findings while carrying no
information about the underlying condition at any point of its range, and the
appendix works a case in which predictions span a factor of \nCalibSpan{}
across deciles while the underlying problem rate is flat in all ten.

The implication is narrow. Reporting that a risk model achieves a given AUC
against recorded findings is a true and useful statement and should be made
in exactly those words. An ordinary train/test split does not detect the
problem, because both halves were produced by the same audit process under
the same targeting policy. What cannot be concluded is that the model has
located where problems are; if the model also influences audit planning, the
validation evidence is partly a record of the model's own effect on the
audit.

\begin{figure}[tbp]
\centering
\includegraphics[width=0.85\textwidth]{ex3_fixed_budget.pdf}
\caption{\capFixedBudget}
\label{fig:budget}
\end{figure}
